\section{Conclusion and Broader Impact}
We aim to support the challenging treatment decisions on ECMO treatment, a scarcely available life-supporting modality for COVID-19 patients. We proposed a disentangled representation model that delivers the propensity score and potential outcomes through semi-supervised variational autoencoding. Several innovative components are proposed and integrated to address the strong selection bias, scarcity of treatment cases, and the curse of dimensionality in patient characterization. The Treatment Joint Inference eliminates the auxiliary networks while regulating factual and counterfactual predictions through input reconstruction, clustering and semi-supervision. The Distribution Balancing disentangles the latent dimensions into different aspects of information and extract the selection bias to aid propensity scoring. The remaining non-confounding dimensions are regulated by maximizing the distributions between treatment and control. The Label Balancing component mitigates data imbalance with generative latent representation while preserving the patient distributions. The experiments on two real-world COVID-19 datasets and the public synthetic dataset show that the model is robust in capturing the underlying decision-making process as well as individual treatment effects, hence outperforming other state-of-the-art algorithms.

\textbf{Potential Impact:}
This work fills the gap between treatment effect models and the critical need for ECMO clinical decision support (or other application domains that face strong selection bias, scarcity of treatment, and overfitting). To our best knowledge, there is no machine learning tool that can aid clinicians in identifying ECMO candidates from high-dimensional EHR data while weighing the risks of disease progression versus the scarcity of ECMO support. 
In fact, the decision is arguably the most complex decision made in the ICU setting. 

The improvement in predictions over state-of-art models leads to the better identification of treatment needs and resource allocation, hence saving lives. When using a machine learning assistive tool in ICU setting, we usually want the model to maximize sensitivity (true positive rate) while fixing a high specificity (true negative rate). Comparing TVAE with existing algorithms (for example, BART), when fixing a high specificity of 0.95, TVAE has a sensitivity of 0.84 (while BART has a sensitivity of 0.71) in the BJC Dataset, and 0.64 (while BART has 0.33) in ISARIC Dataset. Considering the capacity of ECMO resources, such performance improvement leads to 374 more patients (18 in BJC Dataset, and 356 in ISARIC) being correctly identified for ECMO treatment (see more details in Appendix F). Although it is impossible to compare treatment effect estimation without counterfactual outcomes, the improved AUROC and AUPRC in factual outcome prediction suggests potential improvement in characterizing the survival/death impact for each treatment decision.  

\textbf{Limitation:}
This work is not without limitations. In clinical practice, when the treatment decisions are not EHR-related or tracked by the EHR system, the assumption of No Unmeasured Confounding might be violated. Meanwhile, the evaluation of model performance on the counterfactual outcomes remains to be a challenge except using synthetic datasets. Due to the concerns in potential algorithm bias and robustness arose from the data collection, sample size, and underestimation of certain groups, the machine learning predictions should only serve as assistance to clinical decision making.  In health care settings, the treatment effect estimated by our model should be used as only one of the inputs to the clinicians alongside other clinical and ethical considerations. How to incorporate our treatment effect model in clinical decisions should be investigated in future studies.